\documentclass[12pt, a4paper]{article}

% ==========================================
%  Packages Setup
% ==========================================
\usepackage[utf8]{inputenc}
\usepackage{amsmath, amssymb, amsthm}
\usepackage{geometry}
\geometry{top=2.5cm, bottom=2.5cm, left=2.5cm, right=2.5cm}
\usepackage{setspace}
\onehalfspacing
\usepackage{hyperref}
\hypersetup{
    colorlinks=true,
    linkcolor=black,
    urlcolor=blue
}
\usepackage{graphicx}
\usepackage{float}
\usepackage{xcolor}

% ==========================================
%  Theorem Environments
% ==========================================
\newtheoremstyle{break}
  {\topsep}
  {\topsep}
  {\normalfont}
  {}
  {\bfseries}
  {.}
  {\newline}
  {}

\theoremstyle{break}
\newtheorem{theorem}{Theorem}[section]
\newtheorem{definition}{Definition}[section]
\newtheorem{corollary}{Corollary}[theorem]
\renewcommand{\qedsymbol}{$\blacksquare$}

\makeatletter
\renewenvironment{proof}[1][\proofname]{\par
  \pushQED{\qed}%
  \normalfont \topsep6\p@\@plus6\p@\relax
  \trivlist
  \item[\hskip\labelsep\bfseries #1\@addpunct{.}]\mbox{}\newline\ignorespaces
}{%
  \popQED\endtrivlist\@endpefalse
}
\makeatother
\providecommand{\tightlist}{\setlength{\itemsep}{0pt}\setlength{\parskip}{0pt}}

\begin{document}

% ==========================================
%  Title Page
% ==========================================
\begin{titlepage}
    \begin{center}
        \vspace*{2cm}
        {\LARGE \bfseries Lecture note on Ordinary Differential Equation \par}
        \vspace{3cm}
        \begin{tabular}{rl}
            \textbf{Author:} & Windson \\
            \textbf{Course:} & Ordinary Differential Equation \\
        \end{tabular}
        \vspace{1.5cm}

        {\large Department of Applied Mathematics \par}
        {\large National Chung Hsing University \par}
        \vspace{0.3cm}
        145 Xingda Road, South District, Taichung City 40227, Taiwan (R.O.C.) \par
        e-mail: info AT windson.cc

        \vfill
        {\large \textbf{License} \par}
        {\small This document is licensed under the Creative Commons Attribution 4.0 International License (CC BY-SA 4.0).\par}
        \vspace{1cm}
        {\large February 26, 2026 \par}
    \end{center}
\end{titlepage}

% ==========================================
%  Table of Contents
% ==========================================
\pagenumbering{roman}
\tableofcontents
\clearpage
\pagenumbering{arabic}
\setcounter{page}{1}

% ==========================================
%  Main Content
% ==========================================
\section{\texorpdfstring{Integrating Factor \(\mu (x)\)}{2-1 Integrating Factor \textbackslash mu (x)}}\label{integrating-factor-mu-x}

In general \textbf{first-order linear differential equation }in standard form:

\[
  \frac{dy}{dt} + p(t) y = g(t)
\]

In convenient ,we also write the equation in this form:

\[
P(t)\frac{dy}{dt} + Q(t) y = G(t)
\]

then ,to solve this ODE, we tend to find a integrating factor \(\mu (x)\) s.t. \((\mu (t) y)' = k(t)\)

\subsection{\texorpdfstring{How to find the \emph{Integrating Factor }}{How to find the Integrating Factor }}\label{how-to-find-the-integrating-factor}

\[
P(t) \frac{dy}{dt} + Q(t)y = G(t)
\]

\[
\mu (x) P(t) \frac{dy}{dt} + \mu (x) Q(t)y = \mu (t)G(t)
\]

\[
\mu (t) \frac{dy}{dt} + \mu (t) \frac {Q(t)}{P(t)}y = \frac {\mu (t) G(t)}{P(t)}
\]

\[
[\mu (t) y]' = \frac{\mu (t) G(t)}{P(t)}
\]

we can imply that \(\mu '(t) = \frac{Q(t)}{P(t)} \mu (t)\) ,then

\[
\mu (t) = e^{\int_{}^{} \frac{Q(t)}{P(t)} dt}
\]

\begin{center}\rule{0.5\linewidth}{0.5pt}\end{center}

\section{Separable Differential Equation}\label{separable-differential-equation}

The first general first-order differential equation is

\[
\frac{dy}{dt} = f(x ,y)
\]

To identify this class of equation , we write in this form

\[
M(x, y) + N(x ,y)\frac{dy}{dx} = 0
\]

When \(M\) is a function of \(x\) only and \(N\) is a function of \(y\), the equation becomes

\[
M(x)dx+ N(y)dy = 0 
\]

and we can \(\int M(x)dx + N(y)dy = \int 0 dx\)

\begin{center}\rule{0.5\linewidth}{0.5pt}\end{center}

\section{Differences Between Linear and Nonlinear DEs}\label{differences-between-linear-and-nonlinear-des}

\subsection{THM-1 (Existence and Unique for First-Order Linear ODE)}\label{thm-1-existence-and-unique-for-first-order-linear-ode}

\subsubsection{Statement}\label{statement}

If \(p(t)\) is continuous on an open interval \(I\), then \(\exists!\ y=\phi(t)\) satisfying \(y' + p(t)y = q(t)\) for all \(t \in I\), with \(y(t_0)=y_0\).

\subsubsection{Proof}\label{proof}

Since P(t) is continuous on I ,we set \(\mu(t) = e^{\int_{t_0}^{t}p(t)dt} \neq0 \ ,t \in I\)(\(\mu '(t) = \mu(t) p(t)\))

We have \(\mu(t) \in  C'\)(\(\mu \ is \ diff.)\)

Multiplying (1) by \(\mu(t)\), then we obtain \([\mu(t)y]' = \mu(t)q(t)\)

Since q(t) is continuous on I ,,thus

\[
\mu(t) y(t) = \int _{t_0}^{t} \mu(t)q(t) dt +C
\]

\[
y(t) - \mu^{-1}(t) \int_{t_0}^{t} \mu(s)q(s)ds + C
\]

\(y = \phi(t)\) exists , is differential and satisfy (1). Also , from the mitial condition (2), we determine uniquely \(c = y_0\). Hence ,there is only one solution of (1)-(2).

\subsection{THM-2 (First-Order Nonlinear ODE)}\label{thm-2-first-order-nonlinear-ode}

\subsubsection{Statement}\label{statement-1}

Let \(f (t ,y)\)and \(\frac{\delta f}{\delta y}\)be continuous in some rectangle R ,\(\alpha < t <\beta\) ,\(\gamma < y < \delta\)

\(R = \{ (t, y)\  | \ \alpha < t < \beta  ,\ \gamma < y < \delta \}\)

containing \((t_0 ,y_0)\). Then , in some interval \((t_0 -h , t_0 + h) \subset (\alpha ,\beta)\)for some \(h > 0\) ,

\(\exists!\ y = \phi(t)\) such that \(y'(t) = f(t,y)\) and \(y(t_0) = y_0\).

\subsection{Remark}\label{remark}

\begin{itemize}
\item
  \textbf{Thm-2 can imply Thm-1 (linear case)}

  In the case of Thm-1 , \(f(t ,y) = q(t) = p(t)y \in C\) and \(\frac{\delta f}{\delta y}(t ,y) = -p(t) \in C\)

  there is only one solution in some interval \((t_0 +h ,t_0 - h) \subset I\)
\item
  \textbf{Uniqueness}

  We can make sure that the graph of two solution cannot have any intersection.
\end{itemize}

\begin{center}\rule{0.5\linewidth}{0.5pt}\end{center}

\section{Autonomous DEs ,Population Dynamics}\label{autonomous-des-population-dynamics}

\subsection{Autonmous}\label{autonmous}

\subsubsection{Ex1 (Exponential Growth)}\label{ex1-exponential-growth}

\(\phi (t)\)is the popuation of the given species at time \(t\)

\[
\frac{dy}{dt} = f(y)
\]

consider \(\frac{dy}{dt} = ry\)

\begin{itemize}
\tightlist
\item
  r \textgreater{} 0 : increase
\item
  r \textless{} 0 : decline
\end{itemize}

Assume r \textgreater{} 0 ,and \(y(0) = y_0\) ,\(\phi(t) = y_0 e^{rt}\)

\subsubsection{Ex2 (Logistic DEs)}\label{ex2-logistic-des}

\[
r = h(y)
\]

Condiser h(y) s.t. h(y) \(\cong\)r \textgreater{} 0,when y is small

\begin{itemize}
\tightlist
\item
  h(y) decrease as y grows larger
\item
  h(y) \textless{} 0 when y is suff large
\end{itemize}

Assume that \(h(y) = r-ay ,\quad a > 0\)

Consider logistic equation:

\[
\frac{dy}{dt} = (r-ay)y
\]

\[
\frac{dy}{dt} = r(1- \frac{y}{k})y , \ where \ k =\frac{r}{a} 
\]

\begin{itemize}
\tightlist
\item
  Find constant solution:
\end{itemize}

\[
\frac{dy}{dy} = 0 \Longleftrightarrow r(1-\frac{y}{k})y = 0
\]

\[
the\ constant \ solutions \ are \quad y = \phi_1(t) \equiv0 ,\ y = \phi_2 \equiv k
\]

\begin{itemize}
\tightlist
\item
  Study for
\end{itemize}

\[
f( y) = r(1-\frac{y}{k})y = -\frac{r}{k}(y - \frac{k}{2})^2 + \frac{rk}{4}
\]

\[
\frac{dy}{dt} = \begin{cases} > 0 \quad if \ 0 < y < k \\ \quad \\ <0 \quad if \ y > k \\ \end{cases}
\]

By \textbf{Exsistence and Uniqueness Theorem} ,no solution can intersect the equilibrium solution \(y = k\)

\begin{itemize}
\tightlist
\item
  Find \(\frac{d^2y}{dt^2}\).
\end{itemize}

\[
\boxed{\frac{d^2y}{dt^2}}  = \frac{d}{dt}(\frac{dy}{dt}) =\frac{d}{dt}f(y) = f'(y)\frac{dy}{dt} = \boxed{f'(y)f(y)}
\]

\[
\frac{d^2y}{dt^2} = \begin{cases} > 0 \quad if \ 0 < y < \frac{k}{2}\ or \ y > k \\ \quad \\ <0 \quad if \ \frac{k}{2}<y < k \\ \end{cases}
\]

\textbf{Note }: \(y = \frac{k}{2}\)is an inflation points

\begin{center}\rule{0.5\linewidth}{0.5pt}\end{center}

\section{Exact Differential Equations}\label{exact-differential-equations}

\subsection{THM}\label{thm}

Let M ,N ,\(M_y ,N_x\)are continuous in \(R = \{ (x ,y) |\  \alpha < x< \beta ,\ \gamma <y <\delta\)

Then the \(equ(*)\) is exact in \(R\) \(\Longleftrightarrow \ M_y(x ,y) = N_x(x ,y)\)at each point in \(R\)

\subsubsection{proof}\label{proof-1}

\begin{itemize}
\tightlist
\item
  \(\Longrightarrow\)
\end{itemize}

Suppose (*) is an exact differential equation, so there exists \(\psi\) such that \(\frac{\partial \psi}{\partial x} = M(x,y)\) and \(\frac{\partial \psi}{\partial y} = N(x,y)\).

then ,\(M_y(x ,y) = \frac{\partial \psi}{\partial x\partial y} = \frac{\partial \psi}{\partial x \partial y}\), \(M_y\)and \(N_x\)are continuous ,\(\psi_{xy} \ ,\  \psi_{yx}\)are continuous.

Hence ,\(\frac{\partial^2 \psi}{\partial y \partial x} = \frac{\partial^2 \psi}{\partial x \partial y}\) ,\(M_y = N_x\)

\begin{itemize}
\tightlist
\item
  \(\Longleftarrow\)
\end{itemize}

Good Find a \(\psi(x) \ s.t. \ \psi_x = M\ ,\ \psi_y = N\), Integrating \(\psi_x = M \ w.r.t \ x\)

\(\psi(x ,y) = \int_{x_0}^{x}M(s ,y)\ ds + h(y) = Q(x ,y) + h(y)\) ,where \(Q(x ,y) = \int_{x_0}^{x} M(s ,y)\ ds\)and \(x_0\)is some specified constant with \(\alpha < x < \beta\)

Differentiating \(Q(x ,y) + h(y) \quad w.r.t \quad y\)

\(\psi_y(x ,y) = \frac{\partial Q}{\partial y}(x ,y) + h'(y) = N(x ,y)\), Then \(h'(y) = N(x ,y) - \frac{\partial Q}{\partial y}(x ,y)\)

To show \(N(x ,y) - \frac{\partial Q}{\partial y}(x ,y)\)only depends on y

\[
\frac{\partial}{\partial x}(N(x ,y) - \frac{\partial Q}{\partial y}(x ,y)) = N_x - \frac{\partial^2Q}{\partial y \partial x} = N_x - M_y = 0
\]

\begin{center}\rule{0.5\linewidth}{0.5pt}\end{center}

\section{First-Order Difference Equations}\label{first-order-difference-equations}

\subsection{Definition}\label{definition}

\[
y_{n+1} = f(n ,y_n) ,\quad n = 0 ,1 ,2...
\]

It is called first-order difference equation. if f is a function of \(y_n\), the difference equation is linear, otherwise, it is nonlinear. The solutions of difference equation is a sequence of numbers \(y_0 ,y_1 ,y_2\) that satisfy the equation for each n.

\subsection{General case}\label{general-case}

\[
y_{n+1} = f(y_n), \quad n = 0 ,1, 2...
\]

then\(y_1 = f(y_0) ,\quad y_2 = f(y_1) = f(f(y_0) = f^2(y_0)\)

In general , the \(n^{th}\) iterate\(y_n\)is \(y_n = f^n(y_0)\)

\begin{itemize}
\item
  \textbf{Equilibrium solutions}

  Solutions for which \(y_n\)has the same value for all n
\end{itemize}

\subsection{Linear Equations}\label{linear-equations}

\subsubsection{Case-1}\label{case-1}

\[
y_{n+1} = \rho_ny_n,\quad n= 0 ,1 ,2...
\]

In general ,\(y_n = \rho_{n-1}...\rho_0y_0 ,\quad n = 1 ,2 ,...\)

if \(\rho_n = \rho \ \ \forall n \text{ ,the equations becomes } y_n = \rho^n y_0\)

\[
\lim_{n \to \infty}{y_n} = \begin{cases} 0, \quad if \ |\rho| < 1; \\ \quad \\ y_0, \quad if \ \rho = 1; \\ \\ \text{doesn't exist ,  otherwise.} \end{cases}
\]

\subsubsection{Case-2}\label{case-2}

\[
y_{n+1} = \rho y_n + b_n,\quad n = 0,1,2 ...
\]

\[
y_1 = \rho y_0 + b_0
\]

\[
y_2 = \rho(\rho y_0 + b_0) + b_1 = \rho^2 y_0 + \rho b_0 + b_1
\]

\[
y_n = \rho^n y_0 + \sum_{j = 0}^{n-1} \rho^{n-1-j} b_j
\]

In the special case where \(b_n=b\neq 0\) for all \(n\), the difference equation is:

\[
y_{n+1} = \rho y_n + b
\]

\(\text{If }\rho \neq 1, \text{we can write the solution in the more compact form}\)

\[
y_n = \rho^n (y_0 - \frac{b}{1-\rho})+\frac{b}{1 - \rho}
\]

\begin{center}\rule{0.5\linewidth}{0.5pt}\end{center}

\section{Homogeneous DEs with constant coefficients}\label{homogeneous-des-with-constant-coefficients}

\subsection{Definition}\label{definition-1}

Many second-order ordinary differential equations have the form

\[
\frac{d^2y}{dt^2} = f(t ,y ,\frac{dy}{dt})
\]

It is linear if the function has the form

\[
\frac{d^2 y}{dt^2} = g(t)- p(t)\frac{dy}{dt} - q(t) y
\]

In general,

\[
y'' + p(t)y' + q(t)y = g(t)
\]

\[
P(t)y'' + Q(t)y' + R(t)y = 0
\]

if \(g(t) = 0 \Longrightarrow \text{homogeneous}\)

\subsection{Initial Value Problem(IVP)}\label{initial-value-problemivp}

we consider that

\[
\begin{cases}  ay'' + by' +cy = 0\\ \quad \\ y(0) = y_0 ,\quad y'(0) = y_0'\\ \end{cases}
\]

 We start by seeking exponential solutions of the form\(y = e^{rt}\),where r is a parameter to be determined. Then it follows that \(y' = re^{rt}\),\(y' = r^2 e^{rt}\)Substituting these expressions.

\[
(a r^2 + b r + c) e^{rt} = 0
\]

Since \(e^{rt} \neq 0\)

\[
a r^2 + b r  + c = 0
\]

This equation is called the \textbf{characteristic equation }for the differential equation

Assuming that the roots of the characteristic equation are real and different, let them be denoted by \(r_1\) and \(r_2\), where \(r_1 \neq r_2\). Then \(y_1(t) = e^{r_1 t} \ and \ y_2(t) = e ^{r_2 t}\)

Then there exist \(c_1,c_2\in\mathbb{R}\) such that
\[
y = c_1 e^{r_1t} + c_2 e^{r_2t}.
\]

\[
y' = c_1 r_1 e^{r_1 t} + c_2 r_2 e^{r_2 t}
\]

\[
y'' = c_1 r_1^2 e^{r_1 t} + c_2 r_2^2 e^{r_2 t}
\]

Substituting these expression,

\[
ay'' + by' + c = c_1 (ar_1^2 + b r_1 + c)e^{r_1 t} + c_2 (ar_2^2 + b r_2 + c)e^{r_2 t} = 0
\]

Consider the initial value \(y(t_0) = y_0 ,\ y'(t_0) = y_0'\)

\[
\begin{cases}c_1 e^{r_1t_0} + c_2 e ^{r_2t_0}  = y_0\\ \\ c_1 r_1 e^{r_1 t_0} + c_2 r_2 e^{r_2 t_0} = y_0'\\ \end{cases}
\]

\[
c_1 = \frac{y_0'-y_0 r_2}{r_1 - r_2}e^{-r_1t_0},\ c_2 = \frac{y_0r_1 - y_0'}{r_2 - r_1} e^{-r_2 t_0}
\]

\begin{center}\rule{0.5\linewidth}{0.5pt}\end{center}

\section{Solution of Linear Homogeneous Equations; the Wronskian}\label{solution-of-linear-homogeneous-equations-the-wronskian}

\subsection{Differential Operator}\label{differential-operator}

Let p and q be continuous functions on an open interval I , for \(\alpha < t < \beta\). The cases for \(\alpha  = -\infty\)or \(\beta = \infty\)or both , are included.Then for any function \(\phi\) that is twice differential on I , we define the differential operator L by the equation \[
  L[\phi] = \phi'' + p \phi' + q \phi 
  \]

for example ,if \(p(t) = t^2 ,\ q(t) = 1+t,\ \phi(t) = sin3t\)

\[
L[\phi](t) = (sin3t)'' + t^2(sin3t)' + (1+t)sin3t
\]

we usually write this equation in the form

\[
L[y] = y''+p(t)y'+q(t)y = 0
\]

with the initial value \(y(t_0) = y_0,\quad y'(t_0) = y_0'\)

\subsection{Theorem 3.2.1 Existence and uniqueness Theorem}\label{theorem-3.2.1-existence-and-uniqueness-theorem}

Consider the initial value problem

\[
y'' + p(t) y' + q(t) y = g(t),\quad y(t_0) = y_0, \quad y'(t_0) = y_0' 
\]

where p, q ,and g are continuous on an open interval I that contains the point \(t_0\). This problem has exactly one solution \(y = \phi (t)\),and the solution exists throughout the interval I.

\begin{itemize}
\tightlist
\item
  \textbf{Note :} the initial problem has an unique solution on an interval I
\end{itemize}

\subsection{Theorem 3.2.2 Principle of Superposition}\label{theorem-3.2.2-principle-of-superposition}

if \(y_1,\ y_2\)are two solutions of the differential equation

\[
L[y] = y'' + p(t) y' + q(t)y = 0
\]

then the linear combination \(c_1y_1 + c_2 y_2\)is also a solution for any value of the constants \(c_1 ,\ c_2\)

\subsubsection{proof}\label{proof-2}

\[
\begin{aligned}
L[c_1y_1 + c_2 y_2] &= [c_1y_1 + c_2 y_2]'' + p[c_1y_1 + c_2 y_2]' + q[c_1y_1 + c_2 y_2] \\\\
&= c_1y_1'' + c_2y_2'' + c_1py_1' + c_2py_2' + c_1qy_1 + c_2q y_2 \\\\
&= c_1[y_1'' + py_1' + q y_1] + c_2[y_2'' + py_2' + q y_2] \\\\
&= c_1L[y_1] + c_2L[y_2]
\end{aligned}
\]

\subsection{Note}\label{note}

\begin{itemize}
\tightlist
\item
  \(W[y_1,y_2](x)\) is not everywhere zero iff \(c_1 y_1 + c_2 y_2\) contains all solutions of the equation.
\item
  Therefore, \(y = c_1 y_1(t) + c_2 y_2(t)\)is called \textbf{general solution }
\item
  \(y_1\)and \(y_2\) are said to \textbf{form a fundamental set of solutions }iff \(Wy_1 ,y_2 \neq 0\)
\end{itemize}

\begin{center}\rule{0.5\linewidth}{0.5pt}\end{center}

\section{Solution of Linear Homogeneous Equations; the Wronskian}\label{solution-of-linear-homogeneous-equations-the-wronskian-1}

\subsection{Solution of Repeated Roots}\label{solution-of-repeated-roots}

In this section, we are going to solve what if \(r_1 = r_2\)

\[
ay'' + by' +cy = 0 \quad --(a)
\]

if \(r_1 = r_2\), then

\[
y_1(t) = c_1 e^{\frac{-b}{2a}t}
\]

To find second solution, we assume that

\[
y = v(t)y_1(t) = v(t)e^{\frac{-b}{2a}t}
\]

\[
y' = v'(t)e^{\frac{-b}{2a}t} - \frac{b}{2a}v(t)e^{\frac{-b}{2a}t}
\]

\[
y'' = v''(t)e^{\frac{-b}{2a}t} - \frac{b}{a}v'(t)e^{\frac{-b}{2a}t} + \frac{b^2}{4a^2}v(t)e^{\frac{-b}{2a}t}
\]

By substitude equation (a) ,cancelling the factor \(e^{\frac{-b}{2a}t}\), rearranging the items, we find that

\[
v'' = 0
\]

\[
v = c_1 + c_2 t
\]

then we have

\[
y = c_1 e^{\frac{-b}{2a}t} + c_2 t e^{\frac{-b}{2a}t}
\]

Thus \(y\) is a linear combination of the two solutions

\[
y_1(t) = e^{\frac{-b}{2a}t}, \quad y_2(t) = t e^{\frac{-b}{2a}t}
\]

\[
W[y_1 ,y_2](t) = e^{\frac{-b}{2a}}
\]

Since ,\(W[y_1 ,y_2](t)\) never zero, \(y_1\)and \(y_2\)form a fundamental set of solutions.

\subsection{Reduction of Order}\label{reduction-of-order}

Suppose that we know one solution \(y_1(t)\) , not everywhere zero, of

\[
y'' + p(t)y' + q(t)y = 0 \quad --(b)
\]

To find the second solution, let

\[
y = v(t) y_1(t)
\]

\[
y' = v'(t)y_1(t) + v(t)y_1'(t)
\]

\[
y'' = v''(t)y_1(t) + v'(t)y_1'(t) + v'(t)y_1'(t) + v(t)y_1''(t)
\]

Substitude \(y\), \(y'\)and \(y''\)to the equation (b)

\[
y_1 v'' + (2y_1' + py_1)v' + (y_1'' + p(t)y_1' + q(t)y_1)v = 0
\]

Since \(y_1\)is a solution of equation (b).

\[
y_1 v'' + (2y_1' + py_1)v' = 0
\]

Here becomes a first-order differential equation for the function \(v'\)

\begin{center}\rule{0.5\linewidth}{0.5pt}\end{center}

\section{Variation of Parameters}\label{variation-of-parameters}

In this section, we will intoduce the gerneral method in principle at least, it can be applied

to any equation, and it requires no detailed assumptions about the form of the solution.

\subsection{Method of Variation of Parameters}\label{method-of-variation-of-parameters}

In gerneral, we consider

\[
y'' + p(t)y' + q(t)y = g(t)
\]

where \(p\) , \(q\) and \(g\) are continuous functions. We assume that we know

\[
y_h = c_1y_1(t) + c_2 y_2(t)
\]

In this idea, we replace the constant \(c_1\),\(c_2\) by \(u_1(t)\),\(u_2(t)\) \(\text{resp.}\)

\[
y = u_1(t) y_1(t) + u_2(t) y_2(t) \quad --(a)
\]

\[
y' = u_1'(t) y_1(t) + u_1(t)y_1'(t) + u'_2(t) y_2(t) + u_2(t)y_2'(t)
\]

Assume that \(u_1'(t) y_1(t) + u'_2(t) y_2(t) = 0\), we have

\[
y' =  u_1(t)y_1'(t) + u_2(t)y_2'(t)
\]

\[
y'' = u_1'(t)y_1'(t) + u_1(t) y_1''(t) + u_2'(t)y_2'(t) + u_2(t)y''_2(t)
\]

Now, we substitute\(y\),\(y'\)and \(y''\)to the orginal equation. We find that

\[
\begin{aligned}
&u_1(t) \left( y_1''(t) + p(t)y_1'(t) + q(t)y_1(t) \right) \\\\
+&u_2(t) \left( y_2''(t) + p(t)y_2'(t) + q(t)y_2(t) \right) \\\\
+&u_1'(t)y_1'(t) + u_2'(t)y_2'(t) = g(t)
\end{aligned}
\]

where \(y_1\)and \(y_2\)are the homogeneous solutions ,hence the equation above can be reduced to

\[
u_1'(t)y_1'(t) +u_2'(t)y_2'(t)
= g(t) \quad --(b)
\]

Since (a) and (b) forms a system of two linear algebraic equations for derivatives \(u_1'(t)\) and \(u_2'(t)\)of the unknown function. By Carmer's rule, we have:

\[
u_1'(t) = -\frac{y_2(t)g(t)}{W[y_1 ,y_2](t)} ,\quad u_2'(t) = -\frac{y_1(t)g(t)}{W[y_1 ,y_2](t)}
\]

Integrate both of them, and substitude into (a) ,you can get the answer.


\end{document}
